\chapter{Materials and Methods}

The experimental design and salt stress treatments were developed based on established methodologies used in 
previous biochar and salinity studies \textcite{Akhtar2021, Hasanuzzaman2021}.

\section{Experimental site and growth conditions}

\section{Plant material}

\section{Experimental design}

\subsection{Growing media treatments}

Potting was carried out on 15 September 2026. Three growing media treatments were established for each mint species: an unamended peat control, and peat amended with either 15~g or 30~g of biochar per pot. Three replicate pots were prepared for each of the two biochar rates in each species. The biochar was incorporated in a dry state, without prior soaking or pre-wetting.

\textcolor{red}{TODO: Still required for this section -- (i) the pot volume and the mass or volume of peat per pot, without which the 15~g and 30~g rates cannot be expressed as a proportion of the substrate; (ii) the number of replicate pots in the peat-only control; (iii) the biochar feedstock and pyrolysis temperature; and (iv) whether the biochar was mixed through the peat or layered.}

\subsection{Salt stress treatments}

\textcolor{red}{TODO: State the salinity levels applied, the salt used, and how the treatments were imposed and maintained. For reference, published mint salinity work spans a wide range: \textcite{Morshedloo2025} applied 0, 50 and 100~mM NaCl to \textit{Mentha} $\times$ \textit{gracilis} from the six-leaf stage to full flowering, while \textcite{Ounoki2021} used 0, 5, 25 and 50~mM over 14 days and 150~mM over 7 days in spearmint.}

\section{Cultivation practices}

\subsection{Nutrient solution}

Plants were supplied with a modified Hoagland nutrient solution prepared according to \textcite{Hoagland1950}. The composition of the full-strength solution is given in Table~\ref{tab:hoagland}. Iron was supplied as Fe-EDTA rather than the ferric tartrate of the original formulation, and the solution was adjusted to approximately pH 5.5.

\textcolor{red}{TODO: State the strength actually used (full, half, or a progressive schedule), the volume applied per pot, the application interval, and the pH and electrical conductivity to which the solution was adjusted and maintained.}

\begin{table}[htbp]
  \centering
  \caption[Composition of the modified Hoagland nutrient solution]{Composition of the full-strength modified Hoagland nutrient solution \parencite{Hoagland1950}.}
  \label{tab:hoagland}
  \begin{tabular}{llr}
    \toprule
    Component & Salt & Concentration \\
    \midrule
    \multicolumn{3}{l}{\textit{Macronutrients} (mM)} \\
    Potassium nitrate             & KNO$_3$                      & 6.0 \\
    Calcium nitrate               & Ca(NO$_3$)$_2\cdot$4H$_2$O   & 4.0 \\
    Ammonium dihydrogen phosphate & NH$_4$H$_2$PO$_4$            & 1.0 \\
    Magnesium sulfate             & MgSO$_4\cdot$7H$_2$O         & 2.0 \\
    \midrule
    \multicolumn{3}{l}{\textit{Micronutrients} ($\mu$M)} \\
    Boric acid                    & H$_3$BO$_3$                  & 46.0 \\
    Manganese chloride            & MnCl$_2\cdot$4H$_2$O         & 9.1 \\
    Zinc sulfate                  & ZnSO$_4\cdot$7H$_2$O         & 0.77 \\
    Copper sulfate                & CuSO$_4\cdot$5H$_2$O         & 0.32 \\
    Sodium molybdate              & Na$_2$MoO$_4\cdot$2H$_2$O    & 0.11 \\
    Iron (as Fe-EDTA)             & FeSO$_4\cdot$7H$_2$O + EDTA  & 20--40 \\
    \bottomrule
  \end{tabular}
\end{table}

\section{Stress application procedure}

\section{Measured parameters}

\subsection{Growth parameters}

\subsection{Physiological measurements}

\subsection{Laboratory analyses}

\section{Statistical analysis}
